% HD-I19. GPT-6 Astra Pro, /root/imprimitive, inherited effort.
% 21 November 2026. Exact line-bundle formula and componentwise consequences.
\subsection{HD-I19: at most three nonreflection values reduce to one line bundle}
\label{sec:HD-I19}

Continue with the actual imprimitive tower and signature $(n,n)$
of HD-I01, throughout the prompt's range $n\geq5$ and
$m\geq\lceil(n^2-2n+3)/2\rceil$. Let $b$ be the total number
of branch values and let $\nu$ be the number whose monodromies
are not reflections. Write $B=b_1+\cdots+b_\nu$ for these
distinct nonreflection values. Let $Q$ be the reduced divisor on
$Y$ containing one point for every uncolored permutation cycle
over these values, including unramified uncolored cycles.
Put
\begin{equation}
 F=h^*\bigl(T_{\mathbf P^1}(-B)\bigr)\otimes\mathcal O_Y(Q)
       \simeq h^*\mathcal O_{\mathbf P^1}(2-\nu)
                                      \otimes\mathcal O_Y(Q).
                                                        \tag{I19.1}
\end{equation}

For $0\leq\nu\leq3$, the exact forgetful differential on every
individual Hurwitz component has rank
\begin{equation}
 \boxed{\operatorname{rank}d\mathcal F
       =b-\nu-h^0(Y,F)
       =2n+g(Y)-1-h^1(Y,F).}                           \tag{I19.2}
\end{equation}
For $g(Y)\leq n-2$, this is also the exact Weil-period rank at
every point. Thus the remaining constraint in these strata is
the cohomology of the explicitly identified line bundle (I19.1),
with no additional jet matrix rank.

Every component with $\nu\leq1$ has generic Weil-period rank
$2n+g(Y)-1$. For $\nu=1$, the forgetful differential has this
rank everywhere except for the explicitly described subcanonical
case in paragraph~2. Further all-point ranks and actual nonempty
components of ranks $2n$ and $2n-2$ are given below.

\paragraph{1. Exact divisor and kernel.}
At a colored cycle of length $a$, the jet divisor of HD-I16 has
multiplicity $a$; at an uncolored ramified cycle it has
multiplicity $a-1$; an uncolored unramified cycle contributes
zero. The full fibre has multiplicity $a$ in every case.
Consequently there is the exact divisor equality
\begin{equation}
            D_{\rm jet}=h^*B-Q,
       \qquad E(-D_{\rm jet})=F,\quad E=h^*T_{\mathbf P^1}.
                                                        \tag{I19.3}
\end{equation}
This includes the omitted highest uncolored jet: the ramification
centre can move, so it is the coefficient of degree $a-1$ that
is absent, leaving precisely the single missing copy of that point.

The local constraint quotient has a kernel $U$ of dimension
$\nu$. A vector in $U$ gives one common constant target tangent
value over each $b_i$, with all the constrained higher
coefficients zero. Evaluation gives
\[
 H^0(T_{\mathbf P^1})\longrightarrow
                   \bigoplus_{i=1}^{\nu}T_{b_i}\mathbf P^1.
\]
This is surjective for $\nu\leq3$, as follows from
$H^1(\mathbf P^1,\mathcal O(2-\nu))=0$ and the evaluation
sequence. A pulled-back vector field has local expansion
$v(b_i+z^a)=v(b_i)\bmod z^a$, so its restriction gives exactly
these constants for both the colored and uncolored jet lengths.
Hence $U$ lies in the image of $H^0(E)$ in
$H^0(E|_{D_{\rm jet}})$.

The restriction kernel is $H^0(F)$. Therefore the kernel of the
global jet constraint map fits into an exact sequence
\[
 0\longrightarrow H^0(F)\longrightarrow\ker J
       \longrightarrow U\longrightarrow0,
 \qquad \dim\ker J=h^0(F)+\nu.
\]
The tangent fibre of the forgetful map is this kernel modulo
the three target vector fields, by HD-I16. The source dimension
is $b-3$. Subtracting the tangent-fibre dimension proves the
first equality of (I19.2).

The balanced codimension formula gives
$N=2m+2n-b$, while $\deg D_{\rm jet}=N+\nu$. Thus
\[
 \deg F=b-\nu-2n,
 \qquad h^0(F)-h^1(F)=b-\nu-2n+1-g.
\]
Riemann--Roch proves the second equality of (I19.2). For
$g\leq n-2$, the cyclic multiplication map is everywhere
surjective by HD-I12, so its dual is injective. Composition
with the exact forgetful map proves the period assertion.

\paragraph{2. Every one-nonreflection component has full generic rank.}
For $\nu=1$, let $c$ be the fixed-space codimension of its
monodromy. Then $\deg D_{\rm jet}=c\leq m$, and
$\deg F=2m-c\geq m$. For $n\geq6$, the prompt's lower
bound gives $m>2n\geq2g-2$; hence $H^1(F)=0$ everywhere.
For $n=5$, the same conclusion holds for $g\leq5$, since
$m\geq9>2g-2$.

The remaining genus is $g=6$, in which the cyclic cover has
no branch points. Thus every cycle of every branch monodromy
is uncolored. In particular $e=d$ for a generating tuple,
because total color zero on every generator generates only
$G(d,d,m)$. The unique nonreflection has at least one
uncolored cycle, so $c\leq m-1$ and $\deg F\geq m+1$.
This is greater than $2g-2=10$ unless
\[
          n=5,\quad m=9,\quad g=6,\quad c=8.
\]
In this last profile the nonreflection is one uncolored
9-cycle. Let $P$ be its unique ramification point. Its full
fibre is $9P$ and its jet divisor is $8P$, giving the actual
line-bundle identity
\[
                    F\simeq\mathcal O_Y(10P).
\]
Serre duality now gives $h^1(F)=1$ exactly when
$\omega_Y\simeq\mathcal O_Y(10P)$, and gives zero otherwise:
a degree-zero line bundle has a nonzero section exactly when
it is trivial. The forgetful rank is therefore respectively
14 or 15 at these points.

This exceptional condition does not contain a Hurwitz component.
The locus of pointed genus-six curves with canonical divisor
$10P$ has dimension 11. Here is its precise source and map:
the stratum $\mathcal H(10)$ of nonzero holomorphic
differentials has dimension $2g=12$, by Kontsevich--Zorich,
\emph{Connected components of the moduli spaces of Abelian
differentials with prescribed singularities}, Invent. Math.
\textbf{153} (2003), 631--678, equation~(1), p.632.
Their local period-coordinate construction is described in
Section~5.2, pp.662--663. Sending $(Y,\omega)$ to its unique
zero $P$ has fibres exactly the nonzero scalar multiples of
$\omega$, so quotienting by $\mathbf C^*$ gives dimension 11.
The original published article is available from
\href{https://webusers.imj-prg.fr/~anton.zorich/Papers/KontsevichZorich_Inventiones.pdf}
{the author's site}, DOI 10.1007/s00222-003-0303-x.

Over such a pointed curve, a map with unique pole $9P$ is a
function in $H^0(\mathcal O_Y(9P))$ with exact pole order nine.
Riemann--Roch gives
\[
 h^0(9P)=9-6+1+h^0(\omega_Y(-9P))
        =4+h^0(P)=5.
\]
Here $h^0(P)=1$, since a nonconstant function with at most one
simple pole would give a degree-one map to $\mathbf P^1$ and
make the genus zero. After sending the distinguished target
value to infinity, target coordinate changes are exactly the
two-dimensional affine group. Their action on nonconstant
functions is free. Thus the remaining function choices have
dimension three. Simple remaining ramification is an open
condition and adds no parameter.

The unramified cyclic root is finite data. Moreover, the tower
recovers the original Galois cover by normal closure: the core
of $\ker\chi$ in $G$ is trivial, since an element in all its
conjugates fixes each coordinate line with scalar one, and
therefore is the identity. Choices of deck marking are also
finite. The exceptional locus in the Hurwitz space consequently
has dimension at most $11+3=14$. The branch count here is
$b=1+(2m+2n-c)=21$, so every Hurwitz component has dimension
18. None is contained in that exceptional locus.

We have proved that the forgetful map is generically smooth on
every $\nu=1$ component. Its local image contains an open subset
of the relevant cyclic-data component. HD-I12--I14 prove that
the cyclic period map has generic rank $2n+g-1$ on every such
component. Hence every $\nu=1$ component has that same generic
period rank. The $\nu=0$ statement is already HD-I14.

\paragraph{3. Concrete low-degree and negative-degree cases.}
For two nonreflection values, (I19.1) is the actual identity
\[
                       F\simeq\mathcal O_Y(Q).
\]
Write $u=\deg Q$, the total number of uncolored cycles over
these two values. The branch count is $b=2n+u+2$, and
\begin{equation}
       \operatorname{rank}d\mathcal F
                       =2n+u-h^0(\mathcal O_Y(Q)).     \tag{I19.4}
\end{equation}
For $u=0$ this is $2n-1$, recovering HD-I18. For $u=1$ and
$g\geq1$, the degree-one argument above gives $h^0(Q)=1$,
so the forgetful map is everywhere immersive of rank $2n$.
For $u=1,g=0$ its rank is $2n-1$.

For $u=2,g\geq2$, the rank is $2n$ exactly on the locus
where $Q$ is a fibre of an actual degree-two morphism
$Y\to\mathbf P^1$, and is $2n+1$ elsewhere. To verify the
criterion, two independent sections of $\mathcal O_Y(Q)$
have no common base point: a base point would yield a
nonconstant function with at most one simple pole. They
therefore define the stated degree-two morphism. There cannot
be three independent sections: three general sections of this
base-point-free degree-two bundle have no common zero and define
a morphism to $\mathbf P^2$ with nondegenerate image. The
product of its mapping degree and its image degree is two.
A nondegenerate plane curve has degree at least two, so the
image is an irreducible conic and the mapping degree is one.
The resulting finite birational map to the smooth conic is an
isomorphism, contradicting $g\geq1$. Thus $h^0(Q)\leq2$.
Conversely the two pullbacks of target linear sections prove
$h^0(Q)=2$ for such a fibre.
For $g=1$ every degree-two line bundle has two sections, and
for $g=0$ it has three. These conclusions also follow directly
from Riemann--Roch. Whenever $g\leq n-2$, all these ranks
are the exact period ranks as well.

For three nonreflection values,
\[
               F\simeq h^*\mathcal O(-1)\otimes\mathcal O_Y(Q),
               \qquad \deg F=u-m.
\]
Thus $u<m$ gives $H^0(F)=0$ and an everywhere immersive
forgetful map of rank $b-3$. At $u=m$ the rank is $b-4$
when the displayed degree-zero line bundle is trivial, and
$b-3$ otherwise. These are again period ranks in the low
genera. They are direct calculations of (I19.2), rather than
surjectivity assumptions on the jet map.

\paragraph{4. Nonempty rank-$2n$ components for every allowed group.}
Put $a=m-3$. Let $A$ have a long cycle on
$\{1,\ldots,a\}$ of color one, and three fixed cycles with
colors $d-1,u,v$ on indices $m-2,m-1,m$. Choose
\[
\begin{array}{c|c|c|c}
 d&e&u&v\\\hline
 3&1&2&2\\
 3&3&1&2\\
 4&1&2&3\\
 4&2&1&1\\
 4&4&1&3\\
 6&1&2&5\\
 6&2&1&1\\
 6&3&1&2\\
 6&6&1&5
\end{array}
\]
so $u,v$ are nonzero and $u+v\equiv e\pmod d$.
Let $\tau=(1\,m-2)$ be uncolored, and set
$B=A^{-1}\tau$. Multiplication by $\tau$ merges the long
cycle and its complementary fixed cycle into a single
uncolored cycle of length $m-2$. The other two cycles of
$B$ are fixed and colored $d-u,d-v$. Thus $Q$ consists
of exactly one point.

Let $s=(1\,2)$, let $t_1=\tau$, $t_2=(1\,m-1)$,
$t_3=(1\,m)$, and use the same phase-one involutory
reflection $s'$ as in HD-I18. Take
\begin{equation}
 \bigl(A,B,\tau,
     s,s,t_1,t_1,t_2,t_2,t_3,t_3,s',s',
         \underbrace{s,s,\ldots,s,s}_{(n-5)\ \text{pairs}}\bigr).
                                                        \tag{I19.5}
\end{equation}
Its product is one, because $AB\tau=\tau^2=1$ and every
remaining reflection is paired. It has $2n+1$ reflection
entries and two nonreflection entries.

The adjacent conjugates $A^jsA^{-j}$, $0\leq j\leq a-2$,
give the uncolored adjacent transpositions on the long cycle.
Together with $t_1,t_2,t_3$ they generate $\mathfrak S_m$.
The product $s's$ gives a unit diagonal difference, whose
conjugates generate the diagonal augmentation subgroup.
The total color of $A$ generates $e\mathbf Z/d$. This is
the same explicit generation argument as HD-I18, and proves
that (I19.5) generates all of $G(d,e,m)$.

The two nonreflection ramification deficits are $m-4$ and
$m-3$. Together with the $2n+1$ transpositions they give
\[
 2g-2=-2m+(m-4)+(m-3)+(2n+1)=2n-6,
 \qquad g=n-2.
\]
The visible colors are $1,d-1,u,v,d-u,d-v$ and sum to
$3d$. There are six of them, the cyclic cover is connected
by the same primitive line-character argument, and both
Hodge dimensions equal $g-1+3=n$. The product-one generating
tuple supplies actual covers by the cut-sphere construction.
Every resulting component consequently has period rank $2n$
at every point, by (I19.4). For $n=5$ these are rank-ten
components for every allowed group and every allowed $m$.

\paragraph{5. Additional arbitrary profiles with two complete fibres.}
The complete-fibre argument of HD-I18 has a useful extension
beyond $\nu\leq3$. In any branch profile containing two
fixed-vector-free monodromies, imposing just those two fibres
already forces a map variation at fixed cyclic data to be a
pulled-back target vector field. Every such target vector field
satisfies the constraints at all additional branch values.
Therefore the full kernel is still exactly
$h^*H^0(T_{\mathbf P^1})$. The forgetful map is everywhere
immersive of rank $b-3$, regardless of the other branch
profiles. In genera $g\leq n-2$ this is also the everywhere
period rank.

For an explicit new example, start with (I18.5) and replace its
consecutive second $s$ and first $t$ by the single entry $st$.
This is an uncolored three-cycle of fixed-space codimension
two. The first $s$ and last $t$ remain, so the same full-group
generators remain present. The product stays one, both the
total ramification deficit and the visible colors stay the
same, and $g=n-2$ with signature $(n,n)$ still holds.
There are now $b=2n+1$ values and three nonreflections, two
of them fixed-vector-free. These actual nonempty components
have period rank $2n-2$ everywhere. In the line-bundle
description their third branch value contributes $m-2$
uncolored cycles, so $F$ has degree $-2$, in agreement with
paragraph~3.

\paragraph{Audit scope.}
The author read the entire argument, including the local
highest-jet convention, the exact line-bundle identity, the
one-nonreflection borderline dimension bound and its actual
normal-closure map, the full group generation and both new
tuple computations. The published Kontsevich--Zorich dimension
formula and its local-coordinate explanation were reopened
and read at the precise locators above. The separate exact
monomial certificate checks the rank-$2n$ and rank-$(2n-2)$
tuples on finite parameter sets; it does not replace the
universal proofs. Individual components outside the proved
numerical and low-degree cases retain exactly the line-bundle
cohomology gate (I19.2), and for higher genus the explicitly
identified cyclic multiplication kernel from HD-I16.
